\section{Introduction}
\label{sec:introduction}
Knowledge Tracing (KT) estimates a learner's evolving mastery state from sequential interactions and predicts future correctness \citep{corbett1994knowledge,piech2015deep,abdelrahman2023knowledge}. Recent KT models, including memory-based, attention-based, and transformer-style variants, are typically compared using aggregate predictive metrics such as the area under the ROC curve (AUC) and accuracy (ACC) \citep{zhang2017dynamic,ghosh2020context,liu2023simplekt}. These metrics are useful, but they may hide systematic weaknesses on low-frequency knowledge components (KCs), especially when most test interactions are concentrated on dense, frequently observed KCs.

This issue is not merely a reporting detail. In educational decision support, predicted probabilities can be used to decide whether a learner should receive remediation, practice another item, or progress to a new concept. A model that ranks answers well overall may still be poorly calibrated on sparse or limited cold-start KCs. In such cases, high predicted confidence can be misleading even when the model's aggregate AUC appears competitive. Calibration metrics such as Expected Calibration Error (ECE), Brier score, and reliability diagrams therefore provide a complementary diagnostic view \citep{guo2017calibration,brier1950verification,degroot1983comparison}.

This paper treats KT evaluation itself as a reproducibility and diagnostic problem. We do not propose a new KT architecture. Instead, we specify a dataset-agnostic diagnostic protocol that extends standard KT evaluation with train-only KC-frequency stratification, per-stratum calibration diagnostics, concept-level cold-start analysis, and an eight-channel leakage and predictive-sanity audit. The framing is intentionally conservative: the goal is reusable evaluation infrastructure for future KT studies---including graph-based and self-supervised approaches whose gains are expected precisely on sparse concepts---rather than any claim of state-of-the-art performance. The protocol is not only preventative. In our own instantiation, its predictive-sanity channel flagged a prediction--label misalignment in the temporal evaluation pipeline that aggregate metrics alone would not have revealed (Section~\ref{sec:discussion}); we report the audit trail as part of the study.

\subsection{Motivating Diagnostic Gap}
Under standard benchmarking protocols, overall KT performance can be misleading because sparse KCs and limited cold-start KCs expose different predictive and calibration behavior. For example, a model may achieve competitive overall AUC on a learner-based split while showing substantially higher calibration error or unstable AUC estimates on sparse KC strata. This discrepancy is particularly critical in educational decision support, where predicted probabilities are used to make decisions regarding student remediation, progression, or mastery thresholds. If a model's probabilities are poorly calibrated, educational interventions may be triggered inappropriately, despite a deceptively high population-level discrimination score. This motivates dataset-local diagnostics that report performance and reliability at the KC-frequency level, coupled with systematic leakage audits to ensure that evaluation pipelines do not contaminate the results.

In addressing this gap, this paper makes three contributions:
\begin{itemize}
    \item \textbf{Problem 1 (Hidden Stratum Behavior):} Overall AUC/ACC values can mask degradation or metric instability on low-frequency concepts. 
    \textbf{Contribution 1:} We formalize a train-only KC-frequency stratification protocol---with an explicit strict cold-start group ($f_{\mathrm{train}}(c) = 0$) separated from very sparse concepts---that prevents sparse-bucket leakage and enables sample-size-aware KT diagnostics across pre-registered frequency buckets.
    \item \textbf{Problem 2 (Unreliable Probability Estimates):} Aggregate rankings do not reveal whether predicted probabilities correspond to empirical correctness rates, which is crucial for educational decision-making. 
    \textbf{Contribution 2:} We integrate calibration-oriented evaluation---ECE, the Brier decomposition into Uncertainty, Reliability, and Resolution, and reliability diagrams---into bucket-level KT reporting, interpreted jointly with discrimination.
    \item \textbf{Problem 3 (Hidden Evaluation Pipeline Contamination):} KT evaluation pipelines can suffer from subtle, undocumented data leakage and alignment errors across folds or splits. 
    \textbf{Contribution 3:} We establish an eight-channel leakage and predictive-sanity audit checklist, together with a reproducible workflow that exports prediction-level files, diagnostic reports, and paper-ready tables.
\end{itemize}

Our experimental investigation is structured around three core Research Questions (RQs):
\begin{itemize}
    \item \textbf{RQ1 (Stratum-level Predictive Heterogeneity):} Do aggregate evaluation metrics obscure heterogeneous predictive behavior across KC-frequency strata?
    \item \textbf{RQ2 (Calibration Profiles by Frequency):} How do model calibration profiles change across different concept frequency strata?
    \item \textbf{RQ3 (Limited Cold-start Concept Diagnostics):} How do baseline KT models behave on KCs with zero or limited training-fold frequency under temporal splits?
\end{itemize}

RQ1 exercises Contribution 1, RQ2 exercises Contribution 2, and RQ3 draws jointly on Contributions 1 and 3. The remainder of this paper is organized as follows. Section~\ref{sec:related} reviews KT evaluation, sparse-concept analysis, and calibration. Section~\ref{sec:protocol} presents the diagnostic protocol. Section~\ref{sec:experiments} reports the experimental design and diagnostic results. Section~\ref{sec:discussion} discusses implications and limitations.
